\documentclass{article}
% --- watermark (auto-inserted) ---
\usepackage{xcolor}
\usepackage{xfp}% \fpeval; in the kernel since 2022, harmless on older LaTeX
\usepackage{draftwatermark}
\newlength{\watermarkTextWidth}
\AddToHook{begindocument/end}{%
  \SetWatermarkText{\normalfont Generated by The Agentic Scientist}%
  \SetWatermarkFontSize{10pt}%
  \SetWatermarkLightness{0.85}%
  \settowidth{\watermarkTextWidth}{\normalfont\fontsize{10pt}{12pt}\selectfont Generated by The Agentic Scientist}%
  \edef\watermarkScale{\fpeval{0.8 * sqrt(\paperwidth**2 + \paperheight**2) / \watermarkTextWidth}}%
  \edef\watermarkAngle{\fpeval{atand(\paperheight/\paperwidth)}}%
  \SetWatermarkScale{\watermarkScale}%
  \SetWatermarkAngle{\watermarkAngle}%
}
% --- end watermark ---
 % For LaTeX2e
\usepackage{paper,times}

\usepackage{hyperref}
\usepackage{url}
\usepackage{graphicx}
\usepackage{booktabs}
\usepackage{amsmath,amssymb}
\usepackage{subcaption}

\newcommand{\kcap}{k_{\mathrm{cap}}}
\newcommand{\hi}{h}
\newcommand{\lo}{\ell}

\title{How Few Comparisons Flip a Bradley--Terry Leaderboard Pair? Refit-Verified Removal Budgets on Crowd-Vote and Sports Data}
\author{}

\begin{document}
\maketitle

\begin{abstract}
Public leaderboards built from pairwise comparisons fit a Bradley--Terry model and publish a ranking, yet how many comparisons must change before a ranking reverses is rarely measured. We study the smallest set of deleted comparisons, within a cap of 2\% of all rows, that reverses the point-estimate order of one adjacent pair after an exact ridge-regularised Bradley--Terry refit. We measure it on a crowd-vote table (Chatbot Arena, 135,630 comparisons) and an NBA game table (12,900 comparisons). Verified budgets are small: 62 and 4 rows for the top and mid pairs of the crowd table (at most 0.046\% of comparisons) and 62 and 8 rows for the sports table (0.48\% and 0.062\%). Mid-table pairs are cheaper than top pairs, and random deletion flips the top pairs in 0 of 20 draws. However, a static first-order influence ranking with refit bisection reaches every one of these budgets, so the more elaborate search we studied adds no measurable gain, and the crowd-versus-sports contrast is confounded with table size, participants and era. A six-model expert-judgment table did not flip within the cap in development runs, for an unknown reason. All findings are descriptive: four cases, two designs, one rating model.
\end{abstract}

\section{Introduction}
Leaderboards that aggregate ``which of A and B is better'' judgments into Bradley--Terry (BT) ratings are used to choose models, to announce releases, and to compare research systems; Chatbot Arena is a prominent crowd-vote example \citep{https://doi.org/10.48550/arxiv.2403.04132}. The BT model and its many extensions are decades old \citep{Davidson1970OnET,Davidson1969MultivariatePC,Leonard1977AnAB}, and the statistical properties of its estimates under sampling noise are increasingly well understood \citep{Simons1999AsymptoticsWT,Gao2021UncertaintyQI,Li2021ThePO}. What is less often reported is a complementary, \emph{combinatorial} question: how much of the underlying comparison table has to change, whether through noise or deliberate action, before a published rank order reverses?

We make this question operational. Given a comparison table and one adjacent pair of rated participants, the \emph{removal budget} is the smallest number of rows whose deletion makes an exact refit of the same BT objective rank the lower participant above the upper one. Deletion is a deliberately weak perturbation (no inserted or relabelled votes), closely related to the ``dropping a little data'' robustness question of \citet{Broderick2020AnAF}. Because every proposed deletion set is checked by an exact refit, any reported budget is a certificate that a reversal is attainable with that many rows, and therefore an upper bound on the true minimum.

\paragraph{Research question.} Under an exact BT refit (ridge $0.01$, ties as half wins), what is the smallest removal set within a cap of $2\%$ of all rows that reverses an adjacent top-of-table pair or an adjacent mid-table pair, and does this budget differ between a crowd-vote table and a sports-results table when both are measured by one verified protocol?

\paragraph{What we do.} We measure verified removal budgets on two held-out tables of real data: a deduplicated Chatbot Arena crowd-vote table and an NBA game-results table. We compare the studied search configuration, which orders rows by first-order influence on the rating gap and then refines the set with re-linearised greedy batches and a local polish, against three evaluator-owned comparators (random removal, first-order influence only, and opponent-targeted removal) and three validity controls (empty set, label permutation, cross-table copy). A separate pair of development tables, including a six-model expert-judgment table, is used to compare search configurations and is reported separately.

\paragraph{What we find.} (i) Budgets are small in all four held-out cases (4 to 62 of 135,630 crowd rows; 8 to 62 of 12,900 NBA rows), and mid pairs are cheaper than top pairs in both tables. (ii) As a share of comparisons the sports budgets are about 10 and 21 times larger, but the two top-pair budgets are the same number of rows (62), so this contrast largely reflects table size and is confounded with participants and era. (iii) A one-shot influence ranking with refit bisection already attains every budget; the additional search machinery produced no measured improvement. (iv) On the six-model expert-judgment development table, no flip was found within the cap, and our evidence cannot say whether that is a property of the data or of the search.

\paragraph{Contributions.} (1) A precisely specified, evaluator-verified removal-budget protocol with comparators and controls, whose budgets are certificates rather than estimates. (2) Measured budgets for four held-out cases on two evaluation designs, with the crowd-versus-sports contrast and its confounds spelled out. (3) A negative result: first-order influence ordering with exact-refit bisection matches the more elaborate configuration everywhere we measured. (4) An explicit scope statement: the claims concern the declared cases only, not confidence-interval-based rank changes, adversarial insertion or vote rigging, other rating models, or causal effects of evaluation design.

\section{Related Work}
\paragraph{Bradley--Terry models and their estimation.} The BT model has been extended to ties and within-pair order effects \citep{Davidson1970OnET,Davidson1977OnET}, to multivariate and two-dimensional comparisons \citep{Davidson1969MultivariatePC,Causeur2005A2E}, to continuous outcomes and generalised score estimation \citep{Stern2011ModeratedPC,Fageot2023GeneralizedBM}, to covariates and sparse grouped designs \citep{Yan2025InferenceIA,Yan2011GROUPEDSP}, and to Bayesian formulations \citep{Leonard1977AnAB,Hussain2013BayesianIO,Rizvi2026BayesianBF}. Conditions for existence of the maximum likelihood estimate and its consistency are studied by \citet{Butler2004TheEO} and \citet{Bong2021GeneralizedRF}, and ridge-penalised fitting is analysed by \citet{Varin2024TractableRR}. Computation and links to priority derivation and PageRank are treated by \citet{Vojnović2019ConvergenceRO,Sawadogo2019CombiningMA,Genest1999DerivingPF,Selby2024PageRankAT}. We use none of these extensions: our evaluator fixes a plain ridge-penalised BT objective in which a tie counts as half a win, a convention that differs from explicit tie models such as that of \citet{Davidson1970OnET}. The ridge term keeps every refit well defined after rows are removed.

\paragraph{Uncertainty and rank inference.} Sampling-based uncertainty for BT scores and ranks is the subject of \citet{Gao2021UncertaintyQI}, \citet{Makur2024MinimaxHT} (hypothesis testing), \citet{Morel-Balbi2025EstimationOP} (partial rankings from sparse, noisy comparisons) and \citet{Li2021ThePO} (estimation error under general comparison graphs). These works ask how sampling variability limits what a ranking can resolve. We instead fix the observed table and ask how few of its rows determine the order of one adjacent pair. This is a statement about a single realised data set, not about a population, and we report no confidence-interval-based rank events.

\paragraph{Data-dropping robustness.} \citet{Broderick2020AnAF} propose an automatic finite-sample metric for when removing a small fraction of data changes a statistical conclusion, and show it on published applied analyses. Our setting is a specific instance in spirit: the conclusion is the order of one adjacent pair under a BT fit, and the dropped fraction is capped at $2\%$. Unlike an approximation-only metric, we accept a removal set only after an exact refit reverses the pair. \citet{Defresne2025PreferenceEF} use maximum-likelihood preference models for active elicitation, a different goal from auditing a fixed table.

\paragraph{Manipulation of crowd-vote leaderboards.} \citet{https://doi.org/10.48550/arxiv.2501.17858} study vote-rigging strategies against Chatbot Arena and their effect on model rankings. Their setting is adversarial and allows an actor to add votes, whereas ours is a deletion-only perturbation checked by exact refit, with no adversary model, so the two are complementary rather than comparable and we do not compare numbers. Influence-based analysis of single training points is the tool behind our row ordering \citep{https://doi.org/10.48550/arxiv.1703.04730}. Work on other evaluation practices of arena-style platforms, and on MT-Bench-style model-judge evaluation, is relevant background that is not in our verified reference set, and we make no claim of novelty relative to it.

\section{Problem Setup and Method}
\subsection{Bradley--Terry leaderboards and the removal budget}
A comparison table has $N$ rows $(a_i,b_i,y_i)$ over $M$ participants, where $y_i\in\{\text{$a$ wins},\text{$b$ wins},\text{tie}\}$ is encoded as a score $s_i\in\{1,0,\tfrac12\}$. For ratings $\theta\in\mathbb{R}^M$ and $d_i=\theta_{a_i}-\theta_{b_i}$ the fit minimises
\begin{equation}
L(\theta)=\sum_{i=1}^{N}\Big[\log\!\big(1+e^{d_i}\big)-s_i d_i\Big]+\tfrac{\lambda}{2}\lVert\theta\rVert^2,\qquad \lambda=0.01,
\label{eq:bt}
\end{equation}
and participants are ranked by $\theta$. Let $\hi$ be ranked directly above $\lo$. For a removal set $R\subseteq\{1,\dots,N\}$ let $\theta^{(-R)}$ be the exact minimiser of Eq.~\eqref{eq:bt} on the remaining rows. The pair \emph{flips} if $\theta^{(-R)}_{\lo}>\theta^{(-R)}_{\hi}$. With cap $\kcap=\lceil 0.02N\rceil$, the \emph{removal budget} is
\begin{equation}
k^\star=\min\{|R|:\ |R|\le\kcap,\ \text{pair flips under }R\},
\end{equation}
and a search procedure is scored by $1-k/\kcap$ if it returns a flipping set of size $k\le\kcap$ and $0$ otherwise. Because the score is linear in $k$, a score near $1$ means a budget far below the cap. We report budgets in two units: rows, and percent of all $N$ comparisons (derived as $k/N$). The evaluator's ``budget fraction'' metric is $k/\kcap$, a fraction of the allowance and not of the table.

Two kinds of target are used. The \emph{top} pair is ranks 1 and 2 of the full-table fit; the \emph{mid} pair is ranks $\lfloor M/2\rfloor$ and $\lfloor M/2\rfloor+1$. A flip here is a reversal of the \emph{point-estimate} order; no confidence-interval rank event is considered.

\subsection{Why influence ordering should help, and what we test}
Deleting row $i$ perturbs the stationarity condition by $-\nabla\ell_i$, so a Newton step predicts the change of the gap $g=\theta_{\hi}-\theta_{\lo}$ as, in the manner of influence-function analyses of the effect of removing a training point \citep{https://doi.org/10.48550/arxiv.1703.04730},
\begin{equation}
\Delta_i\;\approx\;\frac{(p_i-s_i)\,(v_{a_i}-v_{b_i})}{1-w_i\,q_i},\qquad v=H^{-1}(e_{\hi}-e_{\lo}),
\label{eq:infl}
\end{equation}
where $p_i=\sigma(d_i)$, $w_i=p_i(1-p_i)$, $H$ is the Hessian of Eq.~\eqref{eq:bt} (including the ridge), and $q_i=x_i^\top H^{-1}x_i$ with $x_i=e_{a_i}-e_{b_i}$ gives the leave-one-out leverage correction. Rows with the most negative $\Delta_i$ close the gap fastest to first order. Our use of this approximation differs from a pure influence analysis in that it only orders candidate rows, and every set is then verified by an exact refit. Eq.~\eqref{eq:infl} is a linearisation, so we test two questions: whether an influence order is enough to reach the exact-refit budget, and whether re-linearising and polishing it reduces the budget further.

\subsection{Studied configuration}
The studied configuration (source in the staged evidence for the selected candidate) fits BT by damped Newton iterations with backtracking and an analytic gradient and dense Hessian, and then proceeds in three steps. (a) It ranks rows by Eq.~\eqref{eq:infl} at the full-table fit and finds, by bisection over prefixes of that order with an exact refit at each probe, the smallest prefix of at most $\kcap$ rows that flips the pair. (b) It runs a re-linearised greedy search: after each batch of removals it refits, recomputes Eq.~\eqref{eq:infl} on the remaining rows, and removes a batch of $\lceil 0.6\,n_{\mathrm{lin}}\rceil$ rows, where $n_{\mathrm{lin}}$ is the number of rows the linearised gains need to close the current gap; once the pair flips it bisects over prefixes of the insertion order. (c) It polishes the best set by drop-and-readd local search: it drops the least useful quarter of the incumbent set (halving that fraction after a failure, for at most 14 rounds), refills greedily under a strictly smaller size limit, and keeps any strictly smaller verified set. A set is accepted only if its own exact refit separates $\theta_{\lo}$ above $\theta_{\hi}$ by a margin of $10^{-5}$; the evaluator records its own verified budget for the returned set. Search stops at 60\% of the allotted deadline. Bisection over prefixes assumes that longer prefixes flip once shorter ones do; this is not guaranteed, so the returned $k$ is a verified upper bound on the minimum and minimality is not claimed.

\subsection{Comparators and validity controls}
The evaluator reports three comparators for each case, scored with the same rule as the studied configuration: \textbf{B1}, random removal (20 random orders per case; we report how many flip within the cap); \textbf{B2}, an influence-only removal, which the evaluator exposes as its ``influence reference''; and \textbf{B3}, an opponent-targeted removal. The evaluator source is not part of the staged evidence, so we do not independently verify how B1--B3 are constructed. We rely on their recorded outcomes (flip counts and budgets) and on the names under which they are recorded. Our reading of B2 as one static first-order influence ranking followed by exact-refit prefix bisection, i.e.\ step (a) above without leverage correction, re-linearisation or polishing, is consistent with the inspected source of the starting baseline (Section~\ref{sec:dev}) but is an interpretation, not a verified fact. B3 is described by the evaluator only as an opponent-targeted removal, and we treat it as a hand-built comparator that needs no gradient.

Three controls are recorded alongside the cases: an \emph{empty removal set}, which should not flip any case; a \emph{label-permutation} control on the top pair of each table, for which the evaluator records the fraction of moved labels against its expectation, the rating-spread ratio relative to the real fit, and a non-vacuity indicator; and a \emph{cross-table copy} control, for which it records the number of flips over its trials. We report only these recorded outcomes. The evidence does not state the thresholds or construction details behind the non-vacuity indicator or the copy control, and we make no claim about them.

\subsection{Why these definitions}
The perturbation scale and the cross-dataset comparison are choices, so we state the reasons for them. \emph{Deletion only:} every remaining row is a genuine observation, so a flip cannot be blamed on fabricated data, and an exact refit turns each deletion set into a checkable certificate. Insertion, relabelling and vote rigging would require an adversary model that we do not specify. \emph{Adjacent pairs:} if a removal set puts $\lo$ above $\hi$ and the original ranking placed other participants between them, then at least one adjacent pair in that chain must also have reversed. The smallest budget over adjacent pairs therefore lower-bounds the budget of any non-adjacent reversal, and adjacent pairs are the cheapest rank events. \emph{Top and mid:} the top pair is the position that leaderboards publicise, and the mid pair probes a typical interior position. \emph{A 2\% cap:} it fixes a common tolerance for how much of a table noise or an actor could plausibly alter and makes scores comparable across tables of different size. It is a convention, not an estimate of any real process, and we did not vary it; all held-out budgets sit well below it except the sports top pair, which uses 24\% of its allowance. % @evidence[[derived/final_stats.json#cases.F3.k_over_cap]]

\section{Experimental Protocol}
\subsection{Data}\label{sec:data}
Table~\ref{tab:data} lists the tables. The search procedure receives rows as integer participant indices with a coded outcome; names are not visible to it in the inspected candidate source. The held-out crowd table is the Chatbot Arena 140k preference release (data from the platform of \citealp{https://doi.org/10.48550/arxiv.2403.04132}) after content-hash deduplication against the development tables: from 135,634 raw rows, 2 were removed for content overlap with development data (a further 652 rows are listed in the evaluator metadata as prompt-only overlaps and were kept) and 2 rows of participants with fewer than 100 rows were dropped, leaving 135,630 rows over 52 models. The sports table holds NBA game results from 2006 to 2015 with 12,900 comparisons over 30 teams. The development tables are the Chatbot Arena 55k release (57,477 rows, 64 models) and the MT-Bench human-judgment table (3,355 rows, 6 models), the latter standing in for expert annotation. Dataset commits and licences come from the public task brief and are listed in the Supplement. Model-judge data were not used.

\begin{table}[t]
\centering
\caption{Comparison tables. Caps are $\lceil 0.02N\rceil$ (derived). Held-out tables were never used to choose a search configuration; development tables were.}
\label{tab:data}
\small
\begin{tabular}{llllrrr}
\toprule
Use & Design & Table & Rows $N$ & Participants & Cap $\kcap$\\
\midrule
Held-out & crowd votes & Arena 140k (deduplicated) & 135,630 & 52 models & 2,713\\
Held-out & sports & NBA 2006--2015 & 12,900 & 30 teams & 258\\
Development & crowd votes & Arena 55k & 57,477 & 64 models & 1,150\\
Development & expert judgments & MT-Bench human & 3,355 & 6 models & 68\\
\bottomrule
\end{tabular}
% @evidence[[paper_data.json#final_evaluation.evaluation_evidence.episodes]] held-out rows, models, caps
% @evidence[[derived/task_meta/final_meta.json]] dedup counts
% @evidence[[derived/task_meta/search_meta.json]] development table sizes
\end{table}

\subsection{Cases, seeds and reported uncertainty}
Each table contributes two cases, a top pair and a mid pair, giving four held-out cases (crowd: F1 top, F2 mid; sports: F3 top, F4 mid) plus two permutation controls (F5, F6) on the top pairs. The held-out evaluation ran two seeded episodes (seeds 21 and 22); the development evaluation ran two (seeds 11 and 12). A family score is the mean of the case scores within a design, and the overall score is the mean of the two family scores. Uncertainty is limited by design: there are two episodes, B1 has 20 random draws per case, and the two held-out episodes returned identical per-case budgets. We therefore report no confidence intervals or significance tests, and nothing here is a population-level interval over data sets.

\subsection{Evaluation scopes and qualification}
Search configurations were scored on the development tables, and one frozen configuration was then evaluated in a single sealed held-out evaluation. We keep the two scopes in separate tables and sections and never combine their scores. No search configuration exceeded the starting baseline on the development tables (Section~\ref{sec:dev}), so the frozen selection is not a winner; it is one of three configurations with identical development scores. Its held-out evaluation completed and was recorded as qualified, with a true success indicator. The evidence does not record which gate criteria the qualification decision covers, so we read it only as a statement that the sealed evaluation ran to completion and was recorded as qualified; we do not read it as certifying controls or validity. The control outcomes in Section~\ref{sec:controls} are separate observations taken from the recorded episodes. Qualification is not evidence of superiority over the starting baseline, which was not run on the held-out tables, and it does not extend the findings beyond the declared cases.

\begin{figure}[t]
\centering
\includegraphics[width=\textwidth]{figures/protocol_overview.png}
\caption{Protocol overview. A comparison table is fit with a ridge-penalised Bradley--Terry model; a target adjacent pair (top or mid) is selected; a removal search proposes at most $2\%$ of rows; the evaluator refits exactly and records the verified budget $k$ and score $1-k/\kcap$. Random, influence-only and opponent-targeted removal serve as comparators; empty-set, label-permutation and cross-table copy runs serve as controls. Development tables (Arena 55k, MT-Bench human) and held-out tables (Arena 140k deduplicated, NBA) are kept separate.}
\label{fig:overview}
\end{figure}

\section{Results}
\subsection{Held-out removal budgets}
Table~\ref{tab:cases} and Figure~\ref{fig:main} give the verified budgets of the studied configuration on the four held-out cases. All four flipped within the cap (flip rate 1.0), and every one used a small share of its allowance: the family scores are 0.988 for crowd votes and 0.864 for sports, the overall score is 0.926, and the mean of $k/\kcap$ over the four cases is 0.074. % @evidence[[paper_data.json#final_evaluation.metrics]]
In the crowd table, deleting 62 of 135,630 rows (0.046\%) reverses the top pair and 4 rows (0.003\%) reverse the mid pair. In the sports table, 62 of 12,900 rows (0.48\%) reverse the top pair and 8 rows (0.062\%) reverse the mid pair. % @evidence[[derived/final_stats.json#cases.F1.studied_k]]
The two episodes (seeds 21 and 22) produced identical budgets in every case, so we observe no episode-to-episode variation; with two episodes this is a consistency check on the search procedure, not an estimate of sampling variability. Table~\ref{tab:agg} lists the aggregate held-out metrics next to the comparator references.

\begin{table}[t]
\centering
\caption{Aggregate held-out metrics (budget-saving score $1-k/\kcap$, averaged over cases within a design, then over designs; the four real cases plus controls enter only through the control gap). Reference columns are the evaluator's score for random removal (B1) and influence-only removal (B2) aggregated the same way.}
\label{tab:agg}
\small
\begin{tabular}{lr}
\toprule
Metric & Value\\
\midrule
Overall score (studied configuration) & 0.926\\
\quad crowd family (F1, F2) & 0.988\\
\quad sports family (F3, F4) & 0.864\\
Reference score, random removal (B1) & 0.041\\
Reference score, influence only (B2) & 0.926\\
Flip rate over real cases & 1.0\\
Mean $k/\kcap$ over real cases & 0.074\\
Control gap (permuted minus real, mean score) & $+0.112$\\
\bottomrule
\end{tabular}
% @evidence[[paper_data.json#final_evaluation.metrics]]
\end{table}

\begin{table}[t]
\centering
\caption{Held-out removal budgets. $k$ is the verified number of removed rows; percentages are $k/N$ (derived). B1 reports how many of 20 random orders flipped within the cap; B2 and B3 report the verified budget of each comparator. The studied configuration and B2 coincide in every case.}
\label{tab:cases}
\small
\begin{tabular}{llrrrrrr}
\toprule
Case & Pair & $N$ & $\kcap$ & Studied $k$ (\% of $N$) & B1 flips & B2 $k$ & B3 $k$\\
\midrule
F1 & crowd, top & 135,630 & 2,713 & 62 (0.046\%) & 0/20 & 62 & 650\\
F2 & crowd, mid & 135,630 & 2,713 & 4 (0.003\%) & 6/20 & 4 & 6\\
F3 & NBA, top & 12,900 & 258 & 62 (0.48\%) & 0/20 & 62 & 88\\
F4 & NBA, mid & 12,900 & 258 & 8 (0.062\%) & 0/20 & 8 & 12\\
\bottomrule
\end{tabular}
% @evidence[[derived/final_stats.json#cases]] all cells of this table
\end{table}

\begin{figure}[t]
\centering
\includegraphics[width=\textwidth]{figures/main_budget_by_case.pdf}
\caption{Verified removal budget as a percent of all comparisons (log scale) for the studied configuration, influence-only removal (B2) and opponent-targeted removal (B3), with the 2\% cap shown. Numbers above bars are rows removed. Bars for the studied configuration and B2 coincide in all four cases. Below each category, the number of random orders (of 20) that flipped within the cap.}
\label{fig:main}
\end{figure}

\subsection{Top pairs versus mid pairs}
In both tables the mid pair is cheaper than the top pair: 4 versus 62 rows in the crowd table and 8 versus 62 rows in the sports table. Random removal is informative here: it never flips a top pair (0 of 20 draws, in either table) or the sports mid pair, but flips the crowd mid pair in 6 of 20 draws. So the crowd mid pair is sufficiently close that undirected deletion of up to 2\% of rows sometimes suffices, while an influence-ordered set needs only 4 rows. The registry does not record rating gaps, so we do not test whether smaller gaps explain this ordering.

\subsection{Crowd votes versus sports: a contrast with confounds}\label{sec:contrast}
As a share of comparisons the sports budgets are larger: the top-pair budget is 10.5 times and the mid-pair budget 21.0 times the crowd value (derived from $k/N$). % @evidence[[derived/final_stats.json#ratios]]
The top-pair ratio is exactly the ratio of table sizes, $135{,}630/12{,}900=10.5$, because both top-pair budgets are 62 rows. The mid pairs differ by 4 versus 8 rows. Hence the per-share contrast is largely a denominator effect for the top pairs, and in absolute rows the two designs are equally fragile there. The designs also differ in participants (52 models versus 30 teams), era, outcome granularity, the tie structure, and the match graph. The choice of normaliser also matters. Dividing instead by the mean number of comparisons a participant appears in ($2N/M$, derived: 5,217 for the crowd table and 860 for the sports table) shrinks the contrast but does not remove it: the 62-row top-pair budget is 1.2\% of a crowd participant's comparisons and 7.2\% of a team's (about 6 times), and the mid-pair budgets are 0.077\% and 0.93\% (about 12 times), against 10.5 and 21.0 times under $k/N$. % @evidence[[derived/final_stats.json#cases.F1.studied_percent_of_mean_participant_comparisons]]
Neither normaliser is privileged: the share of the table suits a perturbation of a fixed fraction of all votes, the share of a participant's comparisons is closer to the evidence behind its rating. Both views give the same qualitative picture (small budgets, mid cheaper than top, sports the larger share), differing in the size of the ratio. We cannot separate these factors with two tables, and we therefore do not attribute the contrast to evaluation design. What the data do support is narrower: under one verified protocol, one adjacent top pair in each of two specific tables can be reversed by deleting 62 rows.

\subsection{The studied configuration equals influence-only removal}
The studied configuration's budget equals that of B2 in all four held-out cases and in both permutation controls (Table~\ref{tab:cases}, Supplement Table~\ref{tab:full}), and the overall held-out score 0.926 equals the evaluator's reference score for B2 to the reported precision. % @evidence[[paper_data.json#final_evaluation.metrics]]
So the re-linearisation, the leverage correction and the polish gave no measured reduction on these cases. The same holds on development tables (Section~\ref{sec:dev}). A single static ranking by first-order influence followed by exact-refit bisection is therefore enough to reach every budget we report, and the budgets can be viewed as properties of that simple procedure's output rather than of the more elaborate one. Opponent-targeted removal needs more rows in all four cases: 650 versus 62, 6 versus 4, 88 versus 62 and 12 versus 8. The gap is large only for the crowd top pair (10.5 times), and modest elsewhere (1.4 to 1.5 times). Random removal's mean overall score is 0.041, against 0.926 for B2. % @evidence[[paper_data.json#final_evaluation.metrics]]

We did not run the starting baseline on the held-out tables. By inspection of its source, its first strategy is the same static influence ranking with prefix bisection as B2 and its second strategy runs only when the first finds no flip within the cap; we therefore expect, but have not measured, that it would reproduce B2's budgets here.

\subsection{Controls}\label{sec:controls}
The empty removal set flipped no case, and the cross-table copy control flipped 0 of 2 trials, which is consistent with the budgets not being artefacts of an already-reversed pair or of row position. The permutation controls were recorded as non-vacuous by the evaluator (Figure~\ref{fig:controls}b): the fraction of moved labels was 0.660 against an expected 0.662 for the crowd table and 0.498 against 0.500 for the sports table, and the rating spread fell to 0.088 and 0.160 of the real value. % @evidence[[derived/final_stats.json#permutation_controls]]
The top pair of a permuted table flipped after only 2 rows (crowd) and 10 rows (sports), compared with 62 rows on the real tables (Figure~\ref{fig:controls}a); the mean score difference between permuted and real top pairs is $+0.112$. % @evidence[[paper_data.json#final_evaluation.metrics.control_gap_permuted_minus_real]]
With outcomes shuffled the ratings collapse together and a handful of rows reorders them. This does not show that real ranks are robust in any absolute sense, and the permuted tables are a lower reference, not a model of noise or manipulation.

\begin{figure}[t]
\centering
\includegraphics[width=0.75\textwidth]{figures/controls_permutation.pdf}
\caption{Label-permutation controls on the held-out tables. (a) Verified top-pair budget in rows (log scale) on the real table and on the label-permuted table. (b) Rating spread of the permuted fit relative to the real fit, with the evaluator's non-vacuity limit of 0.5.}
\label{fig:controls}
\end{figure}

\subsection{Development tables: a three-way tie and an unflipped expert table}\label{sec:dev}
Table~\ref{tab:dev} reports the development evaluation of three search configurations: the influence-guided configuration studied above (Configuration 1) and two further configurations that are identified in the evidence only by their titles, a cell-aggregated relaxed removal search (Configuration 2) and a stationarity-LP target-rating search (Configuration 3). All three scored 0.4993, the same as the starting baseline (difference 0.0); all three were scored (3 of 3, none missing). % @evidence[[paper_data.json#downstream_study]]
For the studied configuration, the search-stage evaluator output splits this score into 0.9987 on the crowd family (Arena 55k) and 0.0 on the expert family (MT-Bench human), with an influence-only reference of 0.4993 and a random reference of 0.0695. These sub-scores are quoted from that logged evaluator output and have no entry in the evidence registry, which records only the three total scores and per-episode runtimes (about 1.3 s, 2.0 s and 114 s). We therefore give no breakdown for Configurations 2 and 3. Their source is not part of the staged inspectable evidence, so we make no statement about their algorithms beyond their titles; we mention them only because their total score did not exceed that of the starting baseline.

\begin{table}[t]
\centering
\caption{Development evaluation (Arena 55k and MT-Bench human; seeds 11 and 12). Scores are not comparable with held-out scores, which use different tables. $^\dagger$Quoted from the logged search-stage evaluator output; no registry entry. Dashes: no grounded breakdown.}
\label{tab:dev}
\small
\begin{tabular}{lrrrrr}
\toprule
Configuration & Score & Crowd & Expert & Flip rate & Time / episode (s)\\
\midrule
Configuration 1 (studied) & 0.4993 & 0.9987$^\dagger$ & 0.0$^\dagger$ & 0.5$^\dagger$ & 1.3\\
Configuration 2 & 0.4993 & -- & -- & -- & 2.0\\
Configuration 3 & 0.4993 & -- & -- & -- & 114\\
\midrule
Starting baseline & 0.4993 & \multicolumn{3}{c}{not broken down} & --\\
\bottomrule
\end{tabular}
% @evidence[[paper_data.json#all_programs]] scores and durations of the three configurations
% @evidence[[paper_data.json#downstream_study]] baseline score
\end{table}

The expert family score of 0.0 with a flip rate of 0.5 over the four non-control cases means that neither the top nor the mid pair of the six-model MT-Bench table flipped within a cap of 68 rows, while both crowd cases did. The two crowd cases together needed three rows on average over episodes (derived from the crowd score and the cap of 1,150). % @evidence[[paper_data.json#all_programs]] crowd score 0.998695..., cap from search_meta.json
Whether the expert table is inherently harder to reverse or the search was myopic cannot be determined: no exhaustive search was run on it, and no expert-design result exists on held-out data. The cause is unknown.

\section{Discussion and Limitations}
\paragraph{What the evidence supports.} Within the declared cases, adjacent-pair reversals under an exact BT refit can be certified with very small deletions: at most 62 rows from tables of 12,900 and 135,630 comparisons. The certified budgets are upper bounds. Far smaller sets could exist if a stronger search were available, so the figures should be read as ``at most this many rows'', and the finding that more elaborate search did not reduce them is a statement about the three configurations tried, not about the minimum.

\paragraph{What it does not support.} (i) \emph{No causal claim about evaluation design.} Crowd votes and sports are confounded with table size, participants, era, outcome granularity and tie structure; one table per design means the contrast cannot be separated from table identity. (ii) \emph{No claim about expert or model-judge designs on held-out data.} The only expert-design evidence is a six-model development table with no flip within the cap and an unknown cause; no model-judge data were evaluated. (iii) \emph{No superiority claim.} The studied configuration equals influence-only removal in every case, and the starting baseline was not run on held-out tables. The held-out evaluation was recorded as qualified, but the qualification criteria are not recorded in the evidence, so this outcome neither establishes that the configuration is better nor extends findings beyond four cases. (iv) \emph{No rank-uncertainty or manipulation claim.} A flip here is a reversal of a point-estimate order for one adjacent pair. We examine neither confidence-interval-based rank sets, nor insertion, label flipping or rigging by an adversary, nor other rating models; an adversary can usually do more than delete rows, and a platform's published rank may rest on other statistics. (v) \emph{Limited uncertainty.} Two identical episodes, 20 random draws, and four cases do not support significance statements, and we compute no intervals over data sets.

\paragraph{Threats to validity.} The budget depends on the target pair; we study only the top two and the two middle ranks. The 2\% cap truncates the search space; a reversal that would need more rows is reported as a failure and not as a budget. Permuted tables are a convenient null but not a model of realistic noise. Influence ordering uses a one-step Newton approximation. Ties count as half a win by evaluator convention. The development and held-out tables differ in release and preprocessing, and scores from the two scopes are never combined.

\paragraph{Implications for practice.} A one-shot influence ranking plus bisection over exact refits certifies small reversal budgets on tables of this size in seconds on a CPU. A reported adjacent order in these tables rests on few rows relative to the table; this does not say that the quality difference is small, nor that the order is wrong.

\section{Conclusion}
We measured how few comparisons must be deleted to reverse the point-estimate order of one adjacent pair in a Bradley--Terry leaderboard, using exact-refit verification. On two held-out tables of real data, the verified budgets were 62 and 4 rows for the top and mid pairs of a 135,630-row crowd-vote table and 62 and 8 rows for a 12,900-row sports table. Mid pairs were cheaper than top pairs, and label-permuted tables flipped with far fewer rows than real ones. A static first-order influence ranking matched the more elaborate search everywhere, so we claim a measurement and not a better method. The crowd-versus-sports difference in percent of comparisons is largely table size and is confounded with participants and era. The expert-judgment development table did not flip within the cap, for an unknown reason. Natural next steps are more tables per design, rank-interval events, and exact solvers on small tables.\label{lastbodypage}

\clearpage
\bibliography{references}
\bibliographystyle{paper}

\clearpage
\appendix
\section{Supplementary Information}

\subsection{Full per-case results including controls}
Table~\ref{tab:full} lists all six held-out cases. The two seeded episodes (21 and 22) returned identical values in every cell. Evaluator-recorded durations per seeded episode were 4.45 s and 4.51 s.

\begin{table}[h]
\centering
\caption{All held-out cases. F5 and F6 are label-permutation controls of F1 and F3. Scores are $1-k/\kcap$. B1 score is the mean over 20 random orders, with zero for a non-flipping order. Percent columns are $k/N$ (derived).}
\label{tab:full}
\footnotesize
\setlength{\tabcolsep}{4pt}
\begin{tabular}{lrrrrrrrrrr}
\toprule
Case & $N$ & $\kcap$ & $k$ & \% & Score & B1 flips & B1 score & B2 $k$ & B3 $k$ & B3 score\\
\midrule
F1 & 135,630 & 2,713 & 62 & 0.0457 & 0.977 & 0/20 & 0.000 & 62 & 650 & 0.760\\
F2 & 135,630 & 2,713 & 4 & 0.0029 & 0.999 & 6/20 & 0.164 & 4 & 6 & 0.998\\
F3 & 12,900 & 258 & 62 & 0.481 & 0.760 & 0/20 & 0.000 & 62 & 88 & 0.659\\
F4 & 12,900 & 258 & 8 & 0.062 & 0.969 & 0/20 & 0.000 & 8 & 12 & 0.953\\
F5 & 135,630 & 2,713 & 2 & 0.0015 & 0.999 & 7/20 & 0.251 & 2 & 2 & 0.999\\
F6 & 12,900 & 258 & 10 & 0.078 & 0.961 & 0/20 & 0.000 & 10 & 10 & 0.961\\
\bottomrule
\end{tabular}
% @evidence[[derived/final_stats.json#cases]] all cells of this table
\end{table}

\subsection{Controls and qualification details}
The empty-set control did not flip any case. The cross-table copy control recorded 0 flips in 2 trials; its construction is not documented in the staged evidence. The permutation controls report a moved-label fraction of 0.6601 (expected 0.6615) and a rating-spread ratio of 0.0882 for F5, and 0.4981 (expected 0.5000) and 0.1598 for F6, both recorded as non-vacuous by the evaluator. The mean score difference between permuted and real top pairs is $+0.112$. % @evidence[[derived/final_stats.json#permutation_controls]]
The held-out evaluation completed with outcome qualified, and the recorded success indicator is true. The evidence does not record the gate criteria behind this outcome, so we do not interpret it as certifying that controls passed or that the evaluation is valid beyond the cases reported. % @evidence[[paper_data.json#final_evaluation.outcome]]
Diagnostics volunteered by the candidate during evaluation are marked unverified by the evaluator and were not used for any number in this paper. Every reported budget comes from the evaluator's own exact refit.

\subsection{Development tables: additional detail}
Besides the scores in Table~\ref{tab:dev}, the logged search-stage evaluator output for the studied configuration reports a random-removal reference of 0.0695, an influence-only reference of 0.4993, a mean budget fraction (of the cap, averaged over the four non-control cases) of 0.5007, and a permuted-minus-real control gap of 0.3153. These figures are quoted from that log and are not entries of the evidence registry. % @evidence[[paper_data.json#all_programs]] total score only
The mean budget fraction is close to $0.5$ because the two expert cases that did not flip count as the full cap. Search scores and held-out scores belong to different evaluation scopes and tables; the near-identical magnitudes should not be compared. All three development configurations were scored (3 of 3, 0 unscored), so there is no missing development evidence.

\subsection{Algorithmic specification of the studied configuration}
The Newton fit uses analytic gradient $\nabla L=\sum_i (p_i-s_i)(e_{a_i}-e_{b_i})+\lambda\theta$ and Hessian $H=\sum_i w_i\,x_ix_i^\top+\lambda I$, with step halving until the objective stops increasing (at most 60 iterations, gradient tolerance $10^{-9}$). Eq.~\eqref{eq:infl} floors the leverage denominator $1-w_iq_i$ at $10^{-3}$. Step (b) uses a shrink factor $0.6$ and stops when the removed set reaches the current size limit. Step (c) begins at a drop fraction of $0.25$ and halves it after each failure. All fits are warm-started from the previous solution. The candidate's own acceptance test uses a margin of $10^{-5}$ on its own fit; how the evaluator's fit and acceptance test are implemented is not documented in the staged evidence. The candidate source reads the table as integer participant indices with a coded outcome, so model names are not used by it. Source: the selected candidate's \texttt{src/program.py} in the staged evidence. A further file in that directory (\texttt{datasets\_offline.py}) holds loaders for offline analysis, was not imported by the evaluated code, and was not executed; we make no claim from it.

\subsection{Data provenance and reproduction}
Dataset versions and licences, from the public task brief, are in Table~\ref{tab:prov}. Raw data are not redistributed. The held-out tables are prepared by the evaluator; the evaluator metadata records the content-overlap removals and prompt-only overlaps reported in Section~\ref{sec:data}, and we have no further documentation of its preparation or checking procedure. The figures are built by \texttt{scripts/derive\_final\_stats.py} (derived percentages and ratios from the evaluator's recorded episodes) and \texttt{scripts/build\_figures.py} (Figures~\ref{fig:main} and~\ref{fig:controls}); Figure~\ref{fig:overview} is a schematic with no quantitative content. Tables~\ref{tab:cases} and~\ref{tab:full} are filled from \texttt{derived/final\_stats.json}. The seeds are 11 and 12 (development) and 21 and 22 (held-out); Evaluating one candidate on the held-out tables takes seconds on a CPU (4.5 s per recorded episode here). We were not able to give independent timing for a full reproduction from raw data.

\begin{table}[h]
\centering
\caption{Dataset versions and licences as listed in the public task brief.}
\label{tab:prov}
\footnotesize
\begin{tabular}{llll}
\toprule
Table & Source & Version & Licence\\
\midrule
Arena 55k & \texttt{lmarena-ai/arena-human-preference-55k} & commit \texttt{18c29834} & apache-2.0\\
Arena 140k & \texttt{lmarena-ai/arena-human-preference-140k} & commit \texttt{6322995a} & cc-by-4.0\\
MT-Bench human & \texttt{lmsys/mt\_bench\_human\_judgments} & commit \texttt{f7d2896d} & cc-by-4.0\\
NBA Elo games & FiveThirtyEight \texttt{nba-elo} & commit \texttt{4c1ff5e3} & CC-BY-4.0\\
\bottomrule
\end{tabular}
\end{table}

\subsection{LLM Usage}
Large language models were used throughout this work. The search configurations and their code, the analysis scripts, and the text of this manuscript were produced with an LLM-based research assistant operating inside an automated pipeline that proposed configurations, ran a development evaluation, froze one selection, and then ran a sealed held-out evaluation whose results the manuscript reports. The assistant also drafted the related-work discussion from a small retrieved bibliography. Every figure and table in this paper is derived by script from the evaluator's recorded outputs. No claim in the paper relies on an LLM's unverified recollection of the literature; related work is described only at the level of the cited titles.

\subsection{Extended limitations}
Four held-out cases on two designs, two identical episodes and twenty random draws are the entire basis for the held-out claims. The development side adds one expert table of six participants. We did not run exhaustive or exact optimisation on any table, so we cannot say how far the certified budgets are from the minimum. We did not examine whether the reversed pairs are statistically distinguishable under sampling uncertainty, which would be the natural reading of an adjacent-pair flip for a practitioner. The bibliography available for this submission includes one work on vote rigging and one on influence functions but lacks the MT-Bench paper, classical case-deletion diagnostics and work on arena-platform practices beyond the platform description, so positioning relative to those lines is incomplete.

\end{document}
